\scanpage{116}\mspage{61}
ceux qu'on peut \unsure{décrire|dériver} à partir du foncteur identique, sur une des catégories du diagramme \deleted{\unsure{que}} où aient cours deux, \inserted{voire trois,} notions d'\inserted{$W$-}équivalences différentes. Ce sont, je le rappelle

$1^\circ$) Les \deleted{\unsure{c}} quatre catégories
\[
 \text{\uline{Fibsc}}\,S\unsure{\hookrightarrow|\longrightarrow}\text{\uline{Fib}}\,S\longrightarrow\text{\uline{Liss}}\,S\unsure{\longrightarrow|\text{---}}\mathrm{Cat}_{/S},\quad W_S^d\subset W_S
 \tag*{$(*)$}
\]
pour \unsure{la|les} \deleted{\unsure{éq}} $W$-\unsure{équivalence|équivalences} fine droite d'une part, grossière de l'autre, et dualement
\[
 \text{\uline{Cofibsc}}\,S\longrightarrow\text{\uline{Cofib}}\,\deleted{\unsure{sc}}S\longrightarrow\text{\uline{Prop}}\,S\longrightarrow\mathrm{Cat}_{/S},\quad W_S^g\subset W_S
 \tag*{$(*')$}
\]
pour la $W$-équiv. fine gauche d'une part, la grossière de l'autre\textup{;} et \deleted{\unsure{du même}} \inserted{de plus}

$2^\circ$) $\mathrm{Cat}_{/S}$, qui était \unsure{incluse|inclus} dans les deux suites précédentes, et qui est la seule cat. du diagramme donnant lieu à trois notions de $W$-équivalences différentes.
\diagram{\begin{tikzcd}[ampersand replacement=\&,column sep=small,row sep=0pt]
 W_S^d\arrow[dr,phantom,"{\subset}" sloped] \& {}\\
 {} \& W_S\,.\\
 W_S^g\arrow[ur,phantom,"{\subset}" sloped] \& {}
\end{tikzcd}}

Dans chacun des quatre cas \deleted{$\unsure{C,\Sigma,\Sigma}$} ($C$, munie des deux localiseurs $\Sigma\subset\Sigma'$) de la suite $(*)$, \deleted{pour voir la nature} \deleted{\unsure{du}} le foncteur
\[
 C\Sigma^{-1}\longrightarrow C\Sigma'^{-1}
\]
est essent. ``le même'', à cause des th. de comparaison. Il suffit de regarder ce foncteur dans un seul cas, p. ex. un des deux extrêmes $\text{\uline{Fibsc}}\,S$ ou $\mathrm{Cat}_{/S}$. Même remarque pour la suite $(*')$. Donc finalement il reste deux foncteurs remarquables
% archive sheet number: 115
